\chapter{Baseline: TF-IDF + Regresión Logística multinomial}\label{cap:baseline}

\section{Motivación}

El baseline cumple un doble propósito: (i) servir como punto de
referencia honesto frente al modelo neuronal y (ii) aislar la dificultad
intrínseca de la tarea de la sofisticación del modelo. Un baseline
basado en representaciones léxicas y un clasificador lineal es además
\emph{interpretable}: los pesos del modelo se asocian directamente a
términos del vocabulario y permiten un análisis cualitativo.

\section{Especificación del modelo}

Sea $\phi : \mathcal{X} \to \R^d$ la transformación TF-IDF aplicada a
la concatenación $[c\,\Vert\,e]$ del \emph{claim} y la evidencia. Sobre
$\phi(x)$ se ajusta una regresión logística multinomial con
regularización $\ell_2$:
\[
  \hat\theta = \argmin_{W, b}\;
  \frac{1}{N} \sum_i \ell_{\text{CE}}(y_i, \softmax(W \phi(x_i) + b))
  + \frac{\lambda}{2} \lVert W \rVert_F^2.
\]

\paragraph{Hiperparámetros.}
\begin{itemize}[leftmargin=*,itemsep=0.2em]
  \item Vocabulario: $n$-gramas de palabras de orden $\{1, 2\}$,
        \emph{min\_df} $= 2$, \emph{max\_df} $= 0.95$.
  \item Vocabulario máximo: 200\,000 términos.
  \item Optimizador: \texttt{liblinear} para multiclase \emph{one-vs-rest}
        o \texttt{lbfgs} para multinomial directa.
  \item $\lambda$ seleccionado por validación cruzada estratificada
        $k = 5$ sobre la rejilla
        $\{10^{-3}, 10^{-2}, 10^{-1}, 1, 10\}$.
\end{itemize}

\section{Justificación estadística}

La regresión logística multinomial es el \emph{estimador de máxima
verosimilitud} bajo el modelo categórico
$Y \mid X = x \sim \mathrm{Cat}(\softmax(W x + b))$. La regularización
$\ell_2$ corresponde a un \emph{prior} Gaussiano isótropo
$W \sim \mathcal{N}(0, \lambda^{-1} I)$, equivalente a estimación MAP.
La función objetivo es estrictamente convexa, garantizando óptimo
global único.

La validación cruzada estratificada preserva las proporciones de clase
en cada \emph{fold}, lo que reduce la varianza de las estimaciones de
F1 macro respecto a una validación cruzada simple.

\section{Implementación}

El código se encuentra en
\path{research-stats/research_stats/baseline/train_tfidf_lr.py}. Las
métricas se calculan con \path{research_stats/evaluation/metrics.py} y
se serializan en formato \texttt{parquet} para su análisis posterior.

\section{Resultados sobre validación interna}

% TODO(resultados): completar tras ejecución sobre FEVER.
\begin{table}[h]
  \centering
  \begin{tabular}{lcccc}
    \toprule
    Modelo & Acc. & F1 macro & NLL & ECE \\
    \midrule
    TF-IDF + LR ($\lambda = 1$) & --- & --- & --- & --- \\
    \bottomrule
  \end{tabular}
  \caption{Métricas sobre la partición \emph{dev} (media $\pm$ desv.
  típica sobre tres semillas).}
  \label{tab:baseline}
\end{table}

\section{Análisis de los pesos}

Se reportarán los 20 términos con mayor peso positivo y los 20 con
mayor peso negativo por clase, lo que ilustra qué señales léxicas
explota el modelo. Esta tabla se incluye en el anexo por extensión, y
en la memoria sólo se discute cualitativamente.
